\documentclass[../Main.tex]{subfiles}

\begin{document}

\section{Recherche zum Leveldesign} \label{sec:Leveldesign}
Dieses Kapitel bündelt die recherchierten Designprinzipien thematisch in sechs Überkategorien. Die Prinzipien stammen aus zwei Fachbüchern und vier GDC-Vorträgen, deren Auswahl wir in Abschnitt~\ref{sec:VorgehenLeveldesign} beschreiben. Wir beschränken uns dabei vor allem auf Prinzipien, die die Gegnerplatzierung und den Kampfraum betreffen, da diese im Zentrum unserer Erkenntnisfrage stehen. Sie werden dabei besonders in Bezug auf das Hack-and-Slash-Genre verstanden. Wo möglich wird zusätzlich Bezug auf die Referenzspielanalyse (Kapitel~\ref{sec:Referenzspielanalyse}) genommen. Die unten genannten Designprinzipien variieren von abstrakten Denkmustern über praktisch anwendbare Methoden bis hin zu Gamedesign-Prozessen. Das liegt an der Vielfalt der Quellen und den verschiedenen Absichten der Autoren.

\begin{itemize}[leftmargin=1.5cm]
    \item \textcolor{red}{[Prüfen:] Der Vortrag \enquote{Level Design in a Day} wird nach dem Streichen von Walkthrough und Blockworld in diesem Kapitel nicht mehr zitiert, steht aber in Kapitel~1 und Abschnitt~\ref{sec:VorgehenLeveldesign} als verwendete Quelle. Entweder dort entfernen und hier \enquote{drei GDC-Vorträgen} schreiben oder wieder einen Inhalt daraus verwenden.}
\end{itemize}

\subsection{Kernmechanik-zentriertes Design} \label{sec:Kernmechanik}

\textcolor{red}{Die Kernmechanik bildet die Grundlage für das Spielerlebnis eines Videospiels.} Sie ist das grundlegende Verb, das die Spielhandlung trägt.\autocite[110]{totten2014architectural} Bei \game{Super Mario Bros.} ist dies das Springen, bei \game{The Legend of Zelda} das Erkunden.\textcolor{red}{\autocites[110]{totten2014architectural}{nintendo_supermariobros_1985,nintendo_zelda_1986}} \textcolor{red}{Von diesem Kern aus kommen weitere Aktionen hinzu, ein Spiel kann also durchaus Zweit- oder Drittmechaniken haben.\autocite[110]{totten2014architectural}}

Das Leveldesign wird von den Mechaniken getrieben. Das Level ist demnach das Medium, durch das sich die Mechaniken entfalten.\autocite{gdc_tenprinciples2010} Deshalb ist es sinnvoll, jene Objektklasse, die zur Kernmechanik beiträgt, so breit wie möglich im Spiel zu verankern. Ist die Kernmechanik beispielsweise das Zerstören von Objekten, sollte möglichst vieles von dieser Klasse erben, sodass die Mechanik auf die eine oder andere Weise überall im Spiel auftaucht.\autocite{gdc_movingmountains2020}

\begin{itemize}[leftmargin=1.5cm]
    \item \textcolor{red}{Totten (S.~110) sagt nicht, dass ein Spiel \enquote{vor allem wegen} seiner Kernmechanik gespielt wird, sondern dass sie die Grundlage des Spielerlebnisses bildet. Das Beispiel Mario/Zelda stammt ebenfalls von Totten (S.~110) und braucht deshalb auch seine Fussnote.}
    \item \textcolor{red}{[Prüfen:] Auch die Vererbung von der Klasse der Kernmechanik kommt im dritten Transkript nicht vor, weil es nicht Kays Vortrag ist. Falls ihr den echten Vortrag gesehen habt und es dort vorkommt, bleibt die Fussnote, sonst als eigene Folgerung formulieren (\enquote{Daraus schliessen wir, dass \ldots}).}
\end{itemize}

Für \game{Project Polearm} heisst das, dass die Welt auf schnellen, aggressiven Kampf ausgerichtet sein muss und dass beispielsweise alle Gegner von derselben grundlegenden Klasse erben.\footnote{Wie das bei uns umgesetzt ist, beschreibt der \textit{Project Polearm Technical Report}.} Da sich die Mechaniken im Level entfalten, bemessen wir auch unsere Räume daran, wie viel Bewegungsfreiraum sie dem Spieler für diesen Kampf lassen.

\subsection{Kampf-Encounter-Design} \label{sec:Encounter}
Encounter-Design hat viel mit der \textit{Lesbarkeit} eines Levels oder eines Spielabschnitts zu tun. Ein Kampfraum sollte bereits vom Eingang aus lesbar sein: Ziel, Gegnerpositionen und Deckung müssen erkennbar sein, damit der Spieler eine Strategie fassen kann.\autocite{gdc_horizonquests2018} Je nachdem, wie übersichtlich ein Raum ist, muss ihn ein Spieler ein- oder mehrmals scannen. \textcolor{red}{Bei Multiplayer-Spielen lässt sich die Lesbarkeit nach Andrew Yoder an das Können der Spieler anpassen: Gelegenheitsspieler sollten einen Raum mit einem einzigen Rundblick erfassen können, während erfahrene Spieler Räume mit Verstecken bevorzugen, die man mehrmals scannen muss.}\autocite[122--123]{totten2014architectural} \textcolor{red}{Ein Beispiel für solche Räume ist unserer Erfahrung nach} \game{Counter-Strike},\autocite{valve_counterstrike_2000} \textcolor{red}{dessen Karten viele Ecken und Verstecke haben, die man beim Betreten einzeln absuchen muss. Für Einzelspieler-Spiele wie unseres halten wir dagegen schnell erfassbare Räume für sinnvoll. Das deckt sich} mit unserer Beobachtung in \game{NieR: Automata}, wo die klar lesbare Anfangsaufstellung der Gegner sofort einen Überblick gibt (Frage~A3, Kapitel~\ref{sec:Referenzspielanalyse}). Umgekehrt lässt sich die Lesbarkeit unserer Ansicht nach auch bewusst einschränken: In \game{Ninja Gaiden Sigma Plus} tauchen neue Gegner häufig ausserhalb des Sichtfeldes auf, was die Spielhaftigkeit verbirgt (Frage~A3, Kapitel~\ref{sec:Referenzspielanalyse}).

Verschiedene Pfade ermöglichen es dem Spieler, seinen präferierten Spielstil auszuleben, da ein Pfad beispielsweise aggressives Verhalten und ein anderer das Schleichen vereinfacht.\autocite{gdc_tenprinciples2010} Gute Kampfräume bieten dafür mehrere Routen, Flankiermöglichkeiten und Höhenunterschiede. Milde Barrieren wie eine Umkehr-Treppe zwingen den Spieler zudem zum kurzen Innehalten, bevor er in den Kampf stürmt.\autocite{gdc_horizonquests2018}

Auch die Platzierung der Gegner beeinflusst das Laufverhalten des Spielers: Je nach Spiel meidet der Spieler Gegner oder greift sie bewusst an. So können Gegner selbst zur unsichtbaren Wand oder zum Wegweiser werden.\autocite[199]{totten2014architectural} Eine grosse Zahl oder besonders starke Gegner können den Spieler sogar aktiv in eine gewünschte Richtung drängen. Dieses Prinzip wird \textcolor{red}{nach Matthew Weise} als \textit{Shrinking Fortress} bezeichnet und wirkt gerade in eigentlich grossen Räumen, weil Gegnerschwärme den Raum gefühlt verengen.\autocite[199--200]{totten2014architectural} Das deckt sich mit unserer Beobachtung zu \game{Ninja Gaiden Sigma Plus}, wo die Umzingelung durch mehrere Gegner ein einfaches Rückwärtsausweichen verhindert und den Spielfluss beschleunigt (Frage~D3, Kapitel~\ref{sec:Referenzspielanalyse}).

\begin{itemize}[leftmargin=1.5cm]
    \item \textcolor{red}{Totten (S.~122--123) schreibt nicht, dass geübte Spieler schneller scannen, sondern dass nach Andrew Yoder einfache Räume (ein Rundblick) für Gelegenheitsspieler und verwinkelte Räume (mehrmals scannen) für erfahrene Spieler passen. Deshalb sind auch das Counter-Strike-Beispiel und die Überleitung zu NieR angepasst.}
    \item \textcolor{red}{Den Begriff \textit{Shrinking Fortress} übernimmt Totten (S.~199) von Matthew Weise, das muss erkennbar sein.}
\end{itemize}

\subsection{Räumliche Typologien} \label{sec:Raumtypologien}
Die räumliche Gestaltung lässt sich in \textit{Narrow Space}, \textit{Intimate Space} und \textit{Prospect Space} einteilen. Jeder dieser Typen symbolisiert für den Spieler eine andere Situation oder Emotion: Narrow Spaces stehen für Engpässe, Hinterhalte und Klaustrophobie. Intimate Spaces bilden den sicheren Standardfall, etwa einen Tutorialraum. Prospect Spaces sind offen und machen verwundbar, da Angriffe von überall möglich sind, weshalb sie oft als Bossräume dienen.\autocite[161--168]{totten2014architectural} Eine ähnliche Unterscheidung findet sich bei Rogers, der grosse Räume für grössere Gegnergruppen und enge Räume für Einzelkämpfe empfiehlt.\autocite[245]{rogers2014levelup} In \game{Ninja Gaiden Sigma Plus} konnten wir beobachten, dass die Räume im Verlauf gezielt schmaler wurden, was unserer Ansicht nach einem Übergang von Intimate zu Narrow Space entspricht (Frage~M3, Kapitel~\ref{sec:Referenzspielanalyse}).

\subsection{Spannungskurve, Pacing \& Überraschung} \label{sec:Pacing}
Spannung entsteht nicht nur durch Erzählung oder Musik, sondern bereits durch das Levellayout selbst. Spiele setzen dabei oft abwechselnd auf Phasen hoher Intensität und tiefer Entspannung. Dieses Muster wird als \textit{Rollercoaster-Pacing} bezeichnet und wiederholt sich fraktalähnlich auf der Ebene einzelner Abschnitte wie auch des gesamten Spiels. \textcolor{red}{Taylor warnt allerdings, dass dieses Muster dadurch schnell vorhersehbar wird.}\autocite{gdc_tenprinciples2010} Es deckt sich mit der Wellenstruktur, die wir in \game{Ninja Gaiden Sigma Plus} beobachtet haben, wo die Spannung mit jeder Welle steigt und zwischen den Wellen kurz abfällt (Frage~D1, Kapitel~\ref{sec:Referenzspielanalyse}).

Besonders das Unvorhergesehene reizt die Spieler, denn ein zu regelmässiges Pacing wird schnell langweilig. Der sogenannte \textit{Designer's Curse} beschreibt dabei ein Problem der Entwickler selbst: \textcolor{red}{Sie kennen die üblichen Muster von Levels so gut, dass sie beim Spielen voraussehen, was als Nächstes kommt.} Das Gegenmittel ist, etablierte Muster gezielt zu brechen, etwa wie in \game{Dead Space 2},\autocite{visceral_deadspace2_2011} wo ein Abschnitt \textcolor{red}{rund 15 Minuten} ohne Gegnerkontakt bleibt, bevor die Bedrohung eskaliert. \textcolor{red}{Laut Taylor entstand dieser Bruch allerdings nicht aus Absicht, sondern aus Zwängen der Produktion.}\autocite{gdc_tenprinciples2010}

Übergeordnet sollte beim Gestalten eines Levels immer zuerst vom gewünschten Gefühl ausgegangen werden, danach werden die Mechaniken und Hindernisse ausgewählt und eingebaut, die dieses Gefühl erzeugen.\autocite{gdc_tenprinciples2010}

\subsection{Spielergeführtes Lernen / Scaffolding} \label{sec:Scaffolding}
Das \textit{Scaffolding} beschreibt als Leveldesignprinzip, wie neue Mechaniken in einem Level am besten eingeführt werden.\autocite[114]{totten2014architectural} \textcolor{red}{Jedes Level sollte dabei eine zentrale Mechanik einführen, vorführen oder abwandeln.}\autocite{gdc_tenprinciples2010} Sie wird zunächst eigenständig und gefahrlos gezeigt, idealerweise ganz ohne erklärenden Text, und erst danach mit weiteren Mechaniken verflochten, was als \textit{Braiding} bezeichnet wird.\autocite{gdc_movingmountains2020} Konkrete Muster dafür sind das Zelda-Dungeon-Modell, bei dem \textcolor{red}{jeder Dungeon ein neues Item einführt, jeder Raum darin eine neue Anwendung davon lehrt} und der Bosskampf die Beherrschung prüft, sowie die Bethesda-Schleife aus Lernen, Spielen, Herausforderung und Überraschung.\autocite{gdc_tenprinciples2010}

So soll beispielsweise in \game{Project Polearm} ein Fernkämpfer isoliert eingeführt werden, bevor er mit einem Nahkämpfer kombiniert auftritt. Damit lässt sich der taktische Druck, den wir in \game{Ninja Gaiden Sigma Plus} bei der Kombination verschiedener Gegnertypen beobachtet haben, schrittweise aufbauen (Frage~D2, Kapitel~\ref{sec:Referenzspielanalyse}).

\begin{itemize}[leftmargin=1.5cm]
    \item \textcolor{red}{[Prüfen:] Das dritte Transkript ist nicht der Vortrag von Edward Kay: Der Sprecher dort erwähnt Kays Vortrag nur (\enquote{Ed Key \ldots\ has this really good talk on GDC}). Ob \enquote{zuerst gefahrlos zeigen, idealerweise ohne Text} und der Begriff \textit{Braiding} von Kay stammen, ist damit nicht belegt. Totten (S.~114) stützt das gefahrlose Einführen und das spätere Kombinieren, den Begriff Braiding aber nicht.}
\end{itemize}

\subsection{Prototyping \& Iteration} \label{sec:Prototyping}
Level sollten nicht in einem Durchgang entworfen, sondern schrittweise entwickelt werden. Grundlegende Fragen lassen sich bereits vor jeder 3D-Geometrie klären, etwa mit nicht-digitalen Prototypen auf Papier oder mit Spielfiguren. Danach folgt üblicherweise die \textit{Greybox}, eine einfache Blockgeometrie, mit der Raumgrössen und Distanzen getestet werden, bevor die grafische Ausarbeitung beginnt.\autocite[124--126]{totten2014architectural}

Wie zentral diese iterative Arbeitsweise ist, zeigt Valve am Beispiel von \game{Half-Life 2}\autocite{valve_halflife2_2004}: Die Level durchliefen mehrere vollständige Überarbeitungsdurchgänge, und die Entwickler bezeichnen ebendiese mehrfachen Iterationen rückblickend als einen Schlüssel zum Erfolg des Spiels, da sie deutlich bessere Designentscheidungen ermöglichten.\autocite{jacobson2005postmortem} Für \game{Project Polearm} entspricht dieser Ablauf unserem eigenen Vorgehen, den Kampfabschnitt zuerst als Greybox zu bauen und iterativ zu verfeinern, bevor Grafik und Politur folgen.

\subsection{Erkenntnisse für Project Polearm} \label{sec:ZwischenfazitLD}

Im Folgenden fassen wir kurz zusammen, welche Erkenntnisse aus diesem Kapitel für unsere Erkenntnisfrage und unseren Prototyp wichtig sind. Sie sind vor allem für den Teil der Erkenntnisfrage relevant, in dem es um die Gegnerplatzierung und das Leveldesign geht, während die Analyse der Referenzspiele (Kapitel~\ref{sec:Referenzspielanalyse}) vor allem das Gegnerverhalten betrifft.

Als übergeordneten Grundsatz gehen wir bei der Gestaltung unseres Prototyps vom gewünschten Gefühl aus und gestalten den Kampf danach (Abschnitt~\ref{sec:Pacing}). Für unser Spiel stellen wir uns einen schnellen und aggressiven Kampf vor, in dem sich der Spieler grundsätzlich überlegen fühlt (Abschnitte~\ref{sec:Genredefinition} und~\ref{sec:Vergleich}). Daraus folgt für uns, dass wir Räume, Hindernisse und Gegner so entwerfen müssen, dass sie dieses Gefühl unterstützen oder hervorrufen. Die konkreten Werkzeuge, mit denen wir diese Gefühle beeinflussen wollen, sind die folgenden.

Erstens muss ein Raum schnell lesbar sein, damit sich der Spieler rasch eine Strategie zurechtlegen kann (Abschnitt~\ref{sec:Encounter}).

Zweitens sind die Raumtypen Narrow, Intimate und Prospect Space besonders wichtig (Abschnitt~\ref{sec:Raumtypologien}). Wenn wir sie gezielt mischen, können wir sowohl das Pacing des Spiels als auch die Gefühle der Spieler lenken. Dafür müssen wir uns aber bewusst sein, wie die einzelnen Raumtypen wirken. So muss etwa die Anzahl der Gegner an den Raumtyp angepasst werden, wie es auch Rogers empfiehlt.\autocite[245]{rogers2014levelup} \textcolor{red}{Welchem Raumtyp ein Raum entspricht, bemessen wir dabei am Bewegungsfreiraum des Spielers (Abschnitt~\ref{sec:Kernmechanik}).}

Nicht nur die Grösse der Räume, sondern auch die Gegner selbst können den Spieler dazu zwingen, sich anders zu bewegen. Gegnergruppen können einen Raum sogar gefühlt verengen (\textit{Shrinking Fortress}, Abschnitt~\ref{sec:Encounter}). Damit haben wir unserer Ansicht nach ein weiteres Mittel, um die Bewegungen der Spieler vorherzusehen und abzuschätzen, welche Gefühle ein Raum bei ihnen auslöst (Abschnitt~\ref{sec:Raumtypologien}).

Drittens nutzen wir das Rollercoaster-Pacing, das Anspannung und Erholung abwechselt und sich fraktalähnlich auf mehreren Ebenen gleichzeitig wiederholt (Abschnitt~\ref{sec:Pacing}). Damit ist gemeint, dass sowohl der gesamte Spannungsbogen eines Spiels als auch die einzelnen Kampfabschnitte einem solchen Verlauf folgen. \textcolor{red}{Weil es laut Taylor schnell vorhersehbar wird, wollen wir das Muster an einzelnen Stellen bewusst brechen.}

Viertens sind das Scaffolding- und das Braiding-Prinzip für unseren Prototyp zentral (Abschnitt~\ref{sec:Scaffolding}). Übertragen auf Gegner bedeutet das, eine neue Gegnervariante zuerst einzeln und gefahrlos vorzustellen und sie erst danach nach und nach in Gruppen und Kombinationen ins Spiel einzubinden. Auch den Grundsatz, dass \textcolor{red}{jedes Level eine zentrale Mechanik einführen, vorführen oder abwandeln soll}, möchten wir auf unseren Prototyp übertragen, und zwar auf einzelne Spielabschnitte statt auf ganze Level.

All diese Erkenntnisse fliessen zusammen mit jenen aus der Analyse der Referenzspiele (Abschnitt~\ref{sec:Vergleich}) in unseren Prototyp und in die Beantwortung unserer Erkenntnisfrage ein. Im nächsten Kapitel zeigen wir, wie wir dabei konkret vorgehen.

\end{document}
